\chapter{Knowledge Hooks and Autonomic Knowledge Actuation}
\label{chap:knowledge-hooks-aka}

%%
%% Frame Law (load-bearing, never violated):
%%   attempt → hook → admission/refusal → durable motion → receipt → replay → accounting → promotion
%%   No hook, no consequence.
%%   No receipt, no authority.
%%   No replay, no substrate.
%%   No accounting, no promotion.
%% Source: /Users/sac/truex/docs/MANIFESTO.md lines 67–70, 97–99
%%

\section{The Failure of Model Authority}
\label{sec:hooks-model-authority-failure}

The preceding chapter established that consequential action requires grounding in observed
operational state: $A = \mu(O^{*})$. This chapter examines what stands between the observation
and the action --- the admission structure through which $\mu$ operates. That structure is the
knowledge hook.

Before arriving at the hook, it is necessary to eliminate a persistent confusion that corrupts
engineering decisions: the assumption that a language model's output constitutes runtime
authority. It does not.

A language model is a probability distribution over tokens. Given a context window, it predicts
the next plausible sequence. The output is a compression artifact --- a sample from that
distribution. It carries no typed receipt. It fires no admission predicate. It cannot replay its
own execution. It cannot refuse an unlawful input by name. It has no lifecycle stage
authority.\footnote{Source:
\texttt{/Users/sac/process-intelligence/doctrine/PROCESS\_INTELLIGENCE\_IS\_NOT.md} lines 59--62:
``AI summarization has no lifecycle authority. It cannot emit typed receipts. It cannot refuse
unlawful inputs with named laws. It cannot replay an execution.''}

The failure mode is subtle. An LLM output may \emph{describe} a correct action. It may even
describe the correct hook. But description is not invocation. A narrated boundary crossing is
not a boundary crossing. The recipe is not the meal; the audit trail of the recipe is not the
receipt of the meal having been prepared under lawful conditions.

This matters operationally. Enterprise systems built around AI agents that produce prose
recommendations, JSON patches, or configuration diffs --- and then apply those outputs without
an intervening admission boundary --- have no way to distinguish lawful consequence from
plausible-sounding noise. Both feel like progress. Only one is receipted. Only one is
auditable. Only one can be replayed to confirm that the actual state transition conformed to
the declared process model.

The Chatman Equation written in its receipted form makes the requirement explicit: $R$ proves
$A = \mu(O^{*})$. The receipt $R$ is not generated by the LLM. The receipt is generated by
the hook --- the admission boundary that evaluates the proposed action before authorizing
motion. Without the hook, there is no $R$. Without $R$, there is no authority.\footnote{Source:
\texttt{/Users/sac/truex/docs/MANIFESTO.md} line 62: ``The receipt is not decoration. The receipt
is the authority surface.''}

\section{The Hook as Admission Boundary}
\label{sec:hooks-admission-boundary}

A knowledge hook is a deterministic admission/refusal boundary. It consumes a proposed motion,
evaluates it against admissible process law, emits a decision, and manufactures the first
durable proof that motion occurred or was refused. A hook that produces no decision is not a
hook --- it is narration. A lifecycle transition without a hook firing is not a lawful
transition; it is noise.\footnote{Source:
\texttt{/Users/sac/process-intelligence/phd-thesis/research/knowledge-hooks/KNOWLEDGE\_HOOKS\_AND\_AKA\_DOCTRINE.md}
lines 3--4; cross-verified in
\texttt{/Users/sac/process-intelligence/phd-thesis/ledgers/HOOK\_AKA\_SOURCE\_MAP.yaml} line 216.}

Formally, from the authoritative corpus definition:

\begin{definition}[Knowledge Hook --- Formal Triple]
A knowledge hook is a $(\textit{predicate}, \textit{guard}, \textit{action})$ triple generated
from $\Sigma$. It enforces an invariant $Q$ on every $\Delta$ admitted into $\mu(O)$:

\[
  \text{hook}(p, q, a): \Delta \vDash Q_p \;\Rightarrow\; \mu(O \sqcup \Delta) = \mu(O) \sqcup \mu(\Delta)
\]

Each hook is compiled ahead of time into a branchless kernel that runs within the eight-tick
beat.\footnote{Source: \texttt{/Users/sac/chatmangpt/knhk/yawl.txt} lines 1721--1735, verbatim
definition extracted in
\texttt{/Users/sac/process-intelligence/phd-thesis/research/knowledge-hooks/01\_hook\_definition\_map.md}
lines 36--42.}
\end{definition}

The VKG form of the hook is:

\[
  H = (\Delta,\; C,\; E,\; R,\; \Pi)
\]

where $\Delta$ is the proposed operational delta, $C$ is the condition/guard surface,
$E$ is the permitted effect surface, $R$ is the receipt requirement, and $\Pi$ is the replay
proof requirement. The minimal hook invariant is $\Pi(H(\Delta)) = \textit{stable}$: every
hook must be replay-stable or terminally refused.\footnote{Source:
\texttt{/Users/sac/truex/docs/MANIFESTO.md} lines 228--268, extracted in
\texttt{/Users/sac/process-intelligence/phd-thesis/research/knowledge-hooks/01\_hook\_definition\_map.md}
lines 64--96.}

Admission is evaluated by a compound predicate:

\[
  \text{Accept}(\Delta) \;\Leftrightarrow\;
    \text{Check}_\Sigma \;\wedge\;
    \text{Check}_H \;\wedge\;
    \text{Check}_T \;\wedge\;
    \text{Check}_P \;\wedge\;
    \text{Check}_C \;\wedge\;
    \text{Check}_{\text{Fresh}} \;\wedge\;
    \text{Check}_R
\]

These seven checks cover: type conformance to the ontology ($\Sigma$), guard satisfaction
($H$), transition law ($T$), policy ($P$), capability and epoch binding ($C$),
freshness ($\text{Fresh}$), and receipt lineage ($R$). Admission is binary; partial satisfaction
does not admit.

The eight formal components of a knowledge hook are: (1)~the Attempt --- the proposed $\Delta$;
(2)~the Evidence Field --- the typed observation surface; (3)~the Admission Predicate --- the
compound SPARQL ASK query above; (4)~the Refusal Predicate --- emitting a typed
\texttt{Refusal} artifact when admission fails; (5)~the Motion Boundary --- the permitted
effect surface $E$; (6)~the Receipt Obligation --- every admitted firing must emit a
cryptographic receipt; (7)~the Replay Obligation --- every hook must be replay-stable or
terminally refused; (8)~the Accounting/Promotion Path --- all attempts, refusals, successes,
retries, rollbacks, and promotions must be conserved in an accounting layer.\footnote{Source:
\texttt{/Users/sac/process-intelligence/phd-thesis/research/knowledge-hooks/HOOK\_AKA\_CLAIM\_LEDGER.yaml}
claims CLM-HOOK-002 through CLM-HOOK-009, cross-referenced with
\texttt{/Users/sac/truex/docs/MANIFESTO.md} lines 343--364.}

What a knowledge hook is \emph{not} is equally load-bearing. A hook is not middleware. It is
not callback logic. It is not instrumentation. It is not a webhook. It is not a monitoring
point. It is not a plugin registered to an event bus.\footnote{Source:
\texttt{/Users/sac/truex/docs/MANIFESTO.md} line 219: ``It is not middleware, not callback logic,
not instrumentation, not a webhook, and not a monitoring point.''; confirmed by
\texttt{/Users/sac/chatmangpt/knhk/yawl.txt} line 1722: ``They are neither functions nor
listeners --- they are embedded invariants that bind semantic constraints ($\Sigma$, $Q$)
directly to data movement and execution.''}

These distinctions are not semantic preferences. They are operational. Middleware passes
control through; a hook terminates it on refusal. A callback is invoked by an external
event; a hook evaluates an admission predicate and emits a typed decision artifact. The
laws follow:

\begin{quote}
  No hook, no consequence.\\
  No receipt, no authority.\\
  No replay, no substrate.\\
  No accounting, no promotion.
\end{quote}

\noindent These four laws define the system as one of conserved consequence rather than task
execution or answer production.\footnote{Source: \texttt{/Users/sac/truex/docs/MANIFESTO.md}
lines 67--70, verbatim.}

\section{Autonomic Knowledge Actuation}
\label{sec:aka}

Autonomic Knowledge Actuation (AKA) is the closed-loop discipline of self-managing process
execution. It is the full lifecycle by which knowledge becomes lawful receipted consequence:

\begin{quote}
  knowledge $\to$ closure $\to$ actuation $\to$ receipt $\to$ reconstruction $\to$ improved closure
\end{quote}

AKA is not monitoring. It is not dashboards. It is not alerting.\footnote{Source:
\texttt{/Users/sac/process-intelligence/doctrine/AUTONOMIC\_KNOWLEDGE\_ACTUATION.md} lines 11--12,
verbatim.} These negative boundaries are as important as the positive definition.
Automation is a unidirectional trigger with no receipt obligation; AKA is a receipted closed
loop. Every actuation must emit a receipt. Every repair must replay to prove closure. Every
decommissioning must produce a closure receipt.\footnote{Source:
\texttt{/Users/sac/process-intelligence/doctrine/AUTONOMIC\_KNOWLEDGE\_ACTUATION.md} lines 21--24.}

The AKA system has five stages that map to the MAPE-K reference model \citep{kephart2003vision}
--- Monitor, Analyze, Plan, Execute, Knowledge --- each bounded by $Q$ invariants:

\begin{enumerate}
  \item \textbf{Monitor} --- continuously observes execution conformance against declared process law
  \item \textbf{Analyze} --- identifies root causes of model-log divergence using formal process evidence
  \item \textbf{Plan} --- selects remediation actions within authorized elastic subnets
  \item \textbf{Execute} --- actuates corrections with cryptographic receipt emission
  \item \textbf{Know} --- accumulates a typed knowledge base of lawful and unlawful execution patterns
\end{enumerate}

The mathematical typestate for an AKA transition $s_1 \to s_2$ requires three simultaneous
conditions: (i) transition $t \in T$ is enabled in the current marking $M_1 \to M_2$;
(ii) the target state satisfies LTL safety $s_2 \vDash \Phi_\text{Gov}$; and (iii) there
exists a BLAKE3 lineage proof $\Pi$ such that $\text{VerifyProof}(\Pi, s_1, s_2) = \text{True}$.
The transition function is typed as:

\begin{equation}
  \texttt{transition} : \text{State}(s_1, \text{Proof}(s_1)) \to \text{Transition}(t) \to \text{Option}(\text{State}(s_2, \text{Proof}(s_2)))
\end{equation}

where the output is \texttt{None} --- unrepresentable, a compiler failure --- if the safety
invariants are violated.\footnote{Source:
\texttt{/Users/sac/process-intelligence/doctrine/autonomic-knowledge-actuation.md}, formal
typestate definition, cross-verified in
\texttt{/Users/sac/process-intelligence/phd-thesis/research/knowledge-hooks/02\_aka\_definition\_map.md}
lines 88--97.}

The distinction between knowledge retrieval and knowledge actuation is foundational. Knowledge
retrieval is looking up what you know. Knowledge actuation is making what you know
consequential. Actuation means: receipt (issue a typed proof that the knowledge was applied),
replay (the proof can be verified by replaying the evidence chain), and repair (the knowledge
was applied to fix a non-conforming process). Knowledge that cannot actuate is documentation.
Documentation is PARTIAL. Knowledge that actuates with receipt and replay is
ALIVE.\footnote{Source:
\texttt{/Users/sac/process-intelligence/doctrine/AUTONOMIC\_KNOWLEDGE\_ACTUATION.md},
knowledge retrieval/actuation distinction section.}

MAPE-K embedded as knowledge hooks is the only way to keep the discipline while closing the
loop at machine speed.\footnote{Source: \texttt{/Users/sac/knhk/DOCTRINE\_2027.md} line 171,
verbatim.} The 2027 era is characterized by MAPE-K autonomic hooks running at sub-nanosecond
cadence, with each stage bounded by $Q$ invariants and zero LLM involvement on the hot
path.\footnote{Source: \texttt{/Users/sac/knhk/DOCTRINE\_2027.md} line 29 (era table): ``2027
$\to$ Autonomous evolution | MAPE-K autonomic hooks | Sub-nanosecond decisions.''}

\section{AutoInstinct and Compiled Cognition}
\label{sec:autoinstinct}

AutoInstinct (\texttt{ainst}) is the trace-to-instinct compiler: the compiler layer above
\texttt{ccog}. It learns lawful response policies from proof-backed traces, OCEL worlds,
public ontology profiles, and adversarial JTBD tests, then compiles admitted policies into
deployable field packs.\footnote{Source:
\texttt{/Users/sac/dteam/crates/autoinstinct/src/lib.rs} doc comment, lines 1--8:
``AutoInstinct v30.1.1 --- trace-to-instinct compilation.'' Cross-verified in
\texttt{/Users/sac/process-intelligence/phd-thesis/research/knowledge-hooks/03\_autoinstinct\_lineage\_map.md}
lines 44--48.}

The \texttt{ainst} manufacturing pipeline is:

\begin{quote}
  ontology profile
  $\to$ OCEL worlds
  $\to$ trace corpus
  $\to$ motif discovery
  $\to$ candidate $\mu$ policy
  $\to$ generated JTBD tests
  $\to$ gauntlet
  $\to$ compiled field pack
  $\to$ \texttt{ccog} deployment
\end{quote}

\texttt{ccog} (Compiled Cognition) is the execution runtime, not a chatbot runtime. Its core
formula is $U \to O^{*}_U \to C_U \to A_U \to R_U$, where $C_U$ is the compiled cognition
artifact and $R_U$ is the PROV receipt (proof plus provenance). There is a hard surface
between \texttt{ainst} (Control Plane) and \texttt{ccog} (Runtime Plane), ensuring that
manufacturing is separated from execution and proof.\footnote{Source:
\texttt{/Users/sac/dteam/crates/ccog/src/lib.rs} doc comment; confirmed by
\texttt{/Users/sac/dteam/crates/ccog/docs/end\_to\_end\_jtbd.md} C4 Container View:
``ainst = Control Plane | ccog = Runtime Plane | both = Proof Plane.''}

The product artifact is \texttt{CompiledCcogConfig}: the admitted runtime configuration
produced when a BLAKE3-digested \texttt{FieldPackArtifact} is loaded into the \texttt{ccog}
runtime. It is a compiled, gauntlet-admitted representation of lawful response policy as a
nonlinear graph of COG8 closures.\footnote{Source:
\texttt{/Users/sac/dteam/crates/ccog/src/runtime/mod.rs} line 64:
\texttt{pub struct CompiledCcogConfig<const N: usize, const E: usize>} \{ \texttt{pub pack:
LoadedFieldPack, pub nodes: [Cog8Row; N], pub edges: [Cog8Edge; E], pub mcp\_projections:
MCPProjectionTable} \}.}

The proof artifact is \texttt{EvidenceLedger}: it holds POWL64 route proofs --- replayable,
receipt-bearing records of what the runtime actually did. Conformance checking uses Van der
Aalst alignment to measure whether live execution matched the admitted
\texttt{CompiledCcogConfig} topology, via fitness, precision, generalization, simplicity, and
false closure detection.\footnote{Source:
\texttt{/Users/sac/dteam/crates/ccog/src/runtime/conformance.rs} lines 1--4 (module doc) and
line 21: \texttt{pub struct EvidenceLedger \{ pub traces: Vec<Powl64> \}}.}

The compile-away law follows: LLM output enters the system only as an untrusted candidate
observation, never as compiled policy. The boundary is the BLAKE3-digested
\texttt{FieldPackArtifact}: everything before the digest is probabilistic; everything after is
deterministic, admitted, and authoritative.\footnote{Source:
\texttt{/Users/sac/compiled-cognition-hub/governance/PHILOSOPHY.md} lines 8--9:
``Intelligence stops being a service and becomes a deterministic, zero-dependency property of
the binary. It ceases to be an Oracle, and becomes an Angel --- present at the moment of
action, bounded by law, and instantly verifiable.'' Cross-verified by
\texttt{/Users/sac/dteam/crates/autoinstinct/src/lib.rs} LLM module doc: ``Untrusted output
goes through strict admission before becoming corpus.''}

\section{CONSTRUCT8 and Bounded Motion}
\label{sec:construct8}

Motion across a knowledge hook is bounded by CONSTRUCT8: the bounded constructive delta
operator. Every mutation is bounded to at most eight verifiable triple lanes under receipt and
replay.\footnote{Source:
\texttt{/Users/sac/process-intelligence/phd-thesis/research/knowledge-hooks/KNOWLEDGE\_HOOKS\_AND\_AKA\_DOCTRINE.md}
lines 35--36: ``9 means split. Every mutation is bounded to 8 verifiable triples under receipt
and replay.''}

The \texttt{Construct8Packet} has fixed width:

\begin{lstlisting}[language=Rust,caption={Construct8Packet structure},label={lst:c8packet}]
Construct8Packet {
  epoch: u64,         // Logical clock
  law_ref: u64,       // Which law governs this packet
  subjects: [u32; 8],   // 8 subject IRIs
  predicates: [u32; 8], // 8 predicate IRIs
  objects: [u32; 8],    // 8 object IRIs
  valid_mask: u8,     // Which lanes are filled
  emit_mask: u8,      // Which lanes to emit
  receipt_seed: [u8; 32], // BLAKE3 seed
}
\end{lstlisting}

The kernel exports two deterministic operations:
\texttt{construct8\_admission(packet, gate) $\to$ Result<(), Refusal>} and
\texttt{construct8\_receipt(packet, prev\_receipt) $\to$ Receipt}. The kernel's only promise
is: $\text{hash}(A) = \text{hash}(\mu(O))$.\footnote{Source:
\texttt{/Users/sac/knhk/GENESIS\_CORE\_SPECIFICATION.md}, verbatim specification of
\texttt{construct8\_admission} and \texttt{construct8\_receipt} extracted in
\texttt{/Users/sac/process-intelligence/phd-thesis/research/knowledge-hooks/04\_construct8\_motion\_boundary\_map.md}
lines 57--80.}

A \texttt{Construct8Packet} that exceeds eight active lanes is refused with
\texttt{RefusalReason::Need9} --- a durable, auditable law-enforcement record, not a
recoverable error. Need9 means: decompose, sequence, compose, or add another byte lane. It
does not mean widen to \texttt{u16} or \texttt{Vec}.\footnote{Source:
\texttt{/Users/sac/insa/AGENTS.md} line 17: ``Do not widen Need9 first: Need9 means decompose,
sequence, compose, or add another byte lane. It does not mean widen to u16 or Vec.'' Cross-verified
by \texttt{/Users/sac/knhk/GENESIS\_2030\_DFLSS\_CHARTER.md} line 205: ``Split law: Need257 and
Need9 produce lawful splits, not widened hot payloads.''}

Each packet in a split sequence receives its own receipt; the receipt chain links them.
The no-direct-tool-to-state-mutation rule follows directly: no tool may alter graph state
without passing through \texttt{construct8\_admission}. The consequence cell $\Gamma =
\langle\text{attempt, hook, projection, decision, mailbox, receipt, replay,
accounting}\rangle$ is the atomic unit of motion, not the tool invocation that precedes
it.\footnote{Source: \texttt{/Users/sac/truex/docs/MANIFESTO.md} lines 343--364, extracted in
\texttt{/Users/sac/process-intelligence/phd-thesis/research/knowledge-hooks/01\_hook\_definition\_map.md}
lines 180--193.}

\section{Receipts, Replay, Accounting, and Promotion}
\label{sec:receipts-replay-aka}

The receipt is the authority surface --- not a hash, not a log, not telemetry. Logs are
observation. Receipts are institutional memory.\footnote{Source:
\texttt{/Users/sac/knhk/V30\_1\_1\_MANIFESTO.md} line 15: ``Logs are observation. Receipts are
institutional memory.''; confirmed by
\texttt{/Users/sac/knhk/GENESIS\_CORE\_SPECIFICATION.md}: ``Receipt: Deterministic proof of
consequence (not telemetry).''}

The Rust structure of a receipt is:

\begin{lstlisting}[language=Rust,caption={Receipt structure},label={lst:receipt}]
pub struct Receipt {
  pub receipt_id: u64,
  pub epoch: u64,
  pub input_digest: u64,
  pub output_digest: u64,
  pub previous_receipt_hash: [u8; 32],
  pub receipt_hash: [u8; 32],
  pub law_ref: u64,
  pub ticks_used: u32,
  pub tick_budget: u32,
}
\end{lstlisting}

A receipt proves that consequence was derived from lawful closure, chained to prior receipts,
and bounded by \texttt{law\_ref}.\footnote{Source:
\texttt{/Users/sac/knhk/GENESIS\_CORE\_SPECIFICATION.md}, Receipt struct, extracted in
\texttt{/Users/sac/process-intelligence/phd-thesis/research/knowledge-hooks/HOOK\_AKA\_CLAIM\_LEDGER.yaml}
claim CLM-RECEIPT-001.}

Replay is the mechanism that converts a receipt into a verified conformance artifact. Without
replay, a receipt is a claim. With replay, a receipt is proof. The minimal hook invariant
$\Pi(H(\Delta)) = \textit{stable}$ requires that every admitted hook execution can be
replayed against the original process model and produce the same conforming result.

Accounting conserves all attempts, refusals, successes, retries, rollbacks, and promotions.
A proof without accounting has no authorization for stage transition.\footnote{Source:
\texttt{/Users/sac/truex/docs/MANIFESTO.md} line 70: ``No accounting, no promotion.''} The
GALL test battery verifies this chain: GALL-H (hook fires / does not fire), GALL-R (valid
/ invalid receipt), GALL-RP (replay match / divergence), GALL-S (sabotage refusal ---
inject impossible transition, named law refusal required).\footnote{Source:
\texttt{/Users/sac/process-intelligence/phd-thesis/ledgers/HOOK\_AKA\_SEARCH\_INDEX.md} line 138;
confirmed by
\texttt{/Users/sac/process-intelligence/phd-thesis/research/knowledge-hooks/SEARCH\_AGENT2\_RESEARCH\_CLUSTER.md}
line 121.}

The Cost Per Valid Actuation (CPVA) is the operational metric derived from this chain:
the cost of reaching the state where $R$ proves $A = \mu(O^{*})$ holds --- not the cost
of attempting actuation, but the cost of achieving receipted, replayable, lawful
actuation.\footnote{\textit{Author Thesis:} CPVA is defined in the research cluster synthesis
\texttt{/Users/sac/process-intelligence/phd-thesis/research/knowledge-hooks/02\_aka\_definition\_map.md}
lines 255 and 271. No independent primary doctrine file has been confirmed for this term at
the time of writing. It is offered here as an operational metric consistent with the ALIVE
framework, pending primary source confirmation.}

\section{From AI Assistance to Mechanical Intelligence}
\label{sec:mechanical-intelligence}

The transition from AI assistance to mechanical intelligence is not a change in model
capability. It is a change in what is permitted to carry authority.

In the AI assistance model, authority flows from model output. The system asks; the model
answers; the human decides; the system acts. The receipt chain is: inference result
$\to$ human judgment $\to$ action. The accountability gap is the human judgment step, which
is not receipted, cannot be replayed, and carries no formal conformance to declared process
law.

In the mechanical intelligence model, authority flows from hook admission. The system
proposes $\Delta$; the hook evaluates $\text{Accept}(\Delta)$; if admitted, motion proceeds;
a receipt is emitted; replay can verify conformance. The accountability chain is complete,
continuous, and auditable without human intervention on the hot path.\footnote{\textit{Author
Thesis:} This characterization of the transition from AI assistance to mechanical intelligence
--- as a shift in the authority surface from model output to hook admission --- is the
author's synthesis of the compiled cognition doctrine. It is consistent with the corpus but
is not stated in those exact terms in any single primary source. The supporting architecture
is described in \texttt{/Users/sac/compiled-cognition-hub/governance/PHILOSOPHY.md} lines
8--9 and operationalized across \texttt{/Users/sac/dteam/crates/ccog/} and
\texttt{/Users/sac/dteam/crates/autoinstinct/}.}

The phrase ``mechanical intelligence'' is precise: the system is mechanical in the sense
that its behavior is fully determined by the receipt chain and the process law, not by
probabilistic inference. Intelligence, in this context, is the compiled output of the ainst
manufacturing pipeline --- a deterministic, zero-dependency property of the binary, present
at the moment of action, bounded by law, and instantly verifiable.\footnote{Source:
\texttt{/Users/sac/compiled-cognition-hub/governance/PHILOSOPHY.md} lines 8--9.}

This is the correct compression of the progression: from AI assistance (prediction-mediated
action) to autonomic actuation (hook-mediated receipt) to mechanical intelligence (compiled
policy with no LLM on the hot path). Each step narrows the authority surface to what can be
proven by replay.\footnote{\textit{Author Thesis:} The three-step progression (assistance
$\to$ actuation $\to$ compiled intelligence) is the author's framing. Individual steps are
grounded in corpus evidence; the progression as a named three-stage arc is the author's
synthesis.}

\section{Implications for Capital-Flow Infrastructure}
\label{sec:hooks-capital-flow}

Capital movement requires lawful admission boundaries. A wire transfer authorized by a
language model output is not a lawfully authorized wire transfer; it is a prediction-mediated
instruction with no receipt chain and no replay guarantee.

Knowledge hooks applied to capital-flow infrastructure impose the same admission structure
described in Section~\ref{sec:hooks-admission-boundary}: the proposed capital movement $\Delta$
is evaluated against the compound predicate $\text{Accept}(\Delta)$, covering entity type
conformance, transition law (authorization rules), policy constraints, capability and epoch
binding, freshness (anti-replay), and receipt lineage (the previous settlement receipt
must be present and valid).\footnote{\textit{Author Thesis:} Application of the seven-check
admission predicate to capital movement is the author's extension of the formal hook
structure to the financial settlement domain. The seven checks are source-grounded
(\texttt{/Users/sac/truex/docs/MANIFESTO.md} lines 228--268); the capital-flow application
is the author's derivation.}

Settlement hooks produce BLAKE3-chained receipts binding each authorized movement to a named
law, a prior receipt, and a tick budget --- the same structure as process receipts, applied
to financial ledger entries. DAO admission hooks enforce on-chain governance predicates before
any treasury action is authorized. Post-Quantum Cryptography (PQC) receipt chains ensure that
receipt lineage remains verifiable as classical cryptographic assumptions erode.

In each domain, the frame law holds: no hook, no consequence. Capital that moves without a
hook firing has no lawful receipt. Capital with no lawful receipt has no replayable authority.
Capital with no replayable authority cannot be promoted to a settled state in any audit
framework that treats conformance as a first-class property.
